\chapter{Konklusion og verifikation} \label{chap:konklusion}


Helt fra starten projektperiode, har vi haft fokus på kommunikation. Det, at vi alle i gruppen viste hvad der forventes af os og hvad der skulle laves, var til en fælles forventningsafstemning for det endelig produkt. Da det overordnede design for spillet skulle bestemmes demokratisk, var det vigtigt for os at vi havde et realistisk syn på hvor langt vi kunne nå på de 2 to uger, samt hvor vores evner og talenter lå, så vi nemmere kunne distribuere arbejdet.
Selvom det til tider var op og ned, kan man se på udfyldelse af krav i Diskussions-afsnittet ovenover, at det gik meget godt. Selve strukturen af koden blev hurtigt sat ind i GitHub, som vi brugte som repository til nemt at dele koden med hinanden. Der er altid stor diskussion for hvad der er den bedste måde at stukturer sin kode og filer på. Nemlig et af de punkter vi kunne have gjort bedre var at definitivt blive enige om hvordan vi vil strukturer koden. Strukturen ændres sig meget, da projektet voksede sig mere omfattede, hvor visse filer var dedikere til kun af holde styr på funktioner. Men set i bakspejlet er disse pointer også noget man lære gennem erfaring; at man ved hvordan man vil strukturer sine filer, inden projektet bliver fyldt med funktioner, der afhænger af hinanden. I Fremtiden er det vigtigt at vi slår fast på vores beslutninger, da omrokeringer i planen, selvom disse kan være nemt forudsigelige, kan hurtigt være med at skabe forvirring, især i større, langvarige projekter, i et programmeringssprog som vi forrige for projektet havde ingen til lav erfaring med. Alt i alt har vi skabt et produkt som overholder kravene fra kursets side af, samt anvendt teknikker fra forelæsninger, øvelser og opgaver, som kulminere til en værktøjskasse af fremgangsmetoder, som langsigtet kan bruges igen.
